\documentclass[10pt,a4paper]{amsart}
\usepackage[margin=19mm]{geometry}
\usepackage{amsmath,amssymb,mathtools}
\usepackage[scr=rsfs,cal=euler]{mathalfa}
\usepackage{xfrac}
\usepackage[unicode,hidelinks,bookmarks=false]{hyperref}
\newcommand{\ie}{i.e.\ }
\newcommand{\cf}{cf.\ }
\newcommand{\resp}{respectively\ }
\newcommand{\Subsection}{\subsection}
\providecommand{\nobreakdash}{\nobreak\hbox{-}}
\setlength{\emergencystretch}{2em}
\multlinegap0pt
\usepackage{bm}
\usepackage{bbm}
\usepackage{dsfont}
\usepackage{xspace}
\usepackage{cancel}
\usepackage{mathtools}
\usepackage{amsthm}
\makeatletter
\@ifundefined{theorem}{\newtheorem{theorem}{Theorem}}{}
\@ifundefined{definition}{\newtheorem{definition}[theorem]{Definition}}{}
\@ifundefined{proposition}{\newtheorem{proposition}[theorem]{Proposition}}{}
\makeatother
\title[A general equivalence relation on the global class of morphi]{A general equivalence relation on the global class of morphisms of a category}
\author{Nizar El Idrissi}
\date{Source: arXiv:2602.12363v2 (CC BY 4.0)}
\begin{document}
\maketitle

\noindent\emph{Source and licence.} A general equivalence relation on the global class of morphisms of a category. Version arXiv:2602.12363v2 (CC BY 4.0), distributed under the Creative Commons Attribution 4.0 International licence, as recorded in the arXiv licence metadata for this exact version and shown on the abstract page https://arxiv.org/abs/2602.12363v2. Full rights notice: licenses/arxiv-2602.12363-CC-BY-4.0.txt.\par\smallskip

\noindent\textit{Editorial scope.} Counted unit: the Proposition whose source label is \texttt{proposition-pullback-equivalence-relation-is-the-same-as-the-main-equivalence-relation} in the article (source0.tex lines 444--454, source label \texttt{proposition-pullback-equivalence-relation-is-the-same-as-the-main-equivalence-relation}), delivered together with the article's own complete original proof of it (source0.tex lines 456--547, one \texttt{proof} environment of 92 lines). No mathematics is written, repaired or re-derived: statement and proof are the article's own text line for line. Required context retained uncounted: definition-master (lines 162--181); definition-pullback-equivalence-relation (lines 276--283). Every definition, display and label of those lines is retained unchanged. Omitted from this package and not redistributed: 1--161, 182--275, 284--443, 455--455, 548--884. Nothing inside a retained excerpt is omitted or altered: every retained body is delivered exactly as the article's lines are written, so no omission receipt of any kind accompanies this package. Citations: the retained excerpts carry no citation command at all, so no bibliography entry of the article is transcribed and no reference is redistributed. Reference adaptation: none was needed and none was made; every reference inside the delivered unit resolves inside it, so no unresolved reference or citation is produced. Printed slips kept literal: no printed word, symbol, hypothesis or number of the counted statement or of its proof was altered. Layout and class: the article's own class is \texttt{article} with packages including authblk, inputenc, fontenc, indentfirst, lmodern, mathtools, verbatim, url, booktabs, amsfonts, amsmath, amssymb, amsthm, microtype; the wrapper is the lane's pinned \texttt{amsart} editorial wrapper at 10pt on A4, which restates the theorem-like environments the retained text needs and adds the article's own macro definitions that the retained text needs (none), with extra wrapper packages (bm, bbm, dsfont, xspace, cancel, mathtools). The delivered numbering is the unit's own. Assets: no retained span contains a figure environment, a graphics-inclusion command, a PDF-page inclusion, a table or a TikZ picture; no image file is copied into this package, the package declares no asset dependency and no image file is redistributed with it. Licence: version arXiv:2602.12363v2 of this article is distributed under the Creative Commons Attribution 4.0 International licence, as recorded in the arXiv licence metadata for this exact version and shown on the abstract page https://arxiv.org/abs/2602.12363v2. A full rights notice is delivered at licenses/arxiv-2602.12363-CC-BY-4.0.txt. Characters: every retained excerpt is pure ASCII, so the delivered engine, XeLaTeX with its default Latin Modern text font, reports no missing character.
\begin{definition}
  \label{definition-master}
  Let $\mathcal{C}$ be a category, $\mathcal{D}$ be a 2-category. \\
  Let $\tau_1,\tau_2 : \mathcal{C} \to \mathcal{D}$ be two functors, and let
  \[ \sigma : \begin{cases} \mathsf{Ob(\mathcal{C})} &\to \mathsf{Ob(\mathcal{D})} \\ \mathsf{Hom(c_1,c_2)} &\to \mathsf{Hom(\sigma(c_1),\sigma(c_2))} \end{cases} \]
  be a function (not necessarily a functor) such that $\sigma, \tau_1$, and $\tau_2$ coincide on objects of $\mathcal{C}$. \\
  Further, let
  \[ m : A \to B \]
  and
  \[ \tilde{m} : \tilde{A} \to \tilde{B} \]
  be two morphisms in $\mathcal{C}$.  \\
  We say that $m$ and $\tilde{m}$ are equivalent (with respect to the data $(\mathcal{C},\mathcal{D},\sigma,\tau_1,\tau_2)$), or
  \[ m \simeq \tilde{m} \]
  if there exist 1-morphisms $u_1 : B \to \tilde{B}$, $u_2 : \tilde{A} \to A$, $v_1 : \tilde{B} \to B$, and $v_2 : A \to \tilde{A}$, and four 2-morphisms
  \[ \Phi : \tau_1(u_1) \circ \sigma(m) \circ \tau_2(u_2) \Rightarrow \sigma(\tilde{m}), \]
  \[ \tilde{\Phi} : \sigma(\tilde{m}) \Rightarrow \tau_1(u_1) \circ \sigma(m) \circ \tau_2(u_2), \]
  \[ \Psi : \tau_1(v_1) \circ \sigma(\tilde{m}) \circ \tau_2(v_2) \Rightarrow \sigma(m), \]
  and
  \[ \tilde{\Psi} : \sigma(m) \Rightarrow \tau_1(v_1) \circ \sigma(\tilde{m}) \circ \tau_2(v_2) . \]
\end{definition}

\begin{definition}
  \label{definition-pullback-equivalence-relation}
Let $S$ and $S'$ be sets, where $S'$ is equipped with an equivalence relation $\simeq$. Let $f : S \to S'$ be a function. Define the \textbf{pullback equivalence relation} $\simeq_f$ on $S$ by 
\[
  s \simeq_f \tilde{s} \quad \Longleftrightarrow \quad f(s) \simeq f(\tilde{s}).
\]
If $\simeq$ is seen as a subset of $S' \times S'$, then $\simeq_f$ is the subset of $S \times S$ defined by $\simeq_f = (f^{\triangleq})^{-1}(\simeq)$, where $f^{\triangleq} : S \times S \to S' \times S'$ is the function defined by $f^{\triangleq}(s,\tilde{s}) = (f(s), f(\tilde{s}))$.
\end{definition}

\begin{proposition}
  \label{proposition-pullback-equivalence-relation-is-the-same-as-the-main-equivalence-relation}
  Let $s, \tilde{s} \in S$. Then
  \begin{align*}
    s \;\simeq_f\; \tilde{s} \text{ in } S \text{ (in the sense of definition \ref{definition-pullback-equivalence-relation})} &\iff  m_{s\tilde{s}} \;\simeq\; m_{s\tilde{s}}
    \quad \text{ with respect to the data } \\
    &(\mathrm{Ind}(S), \mathrm{Ind}(S'), \mathsf{\sigma}, \mathsf{\sigma}, \mathsf{\sigma}) \text{ (in the sense of definition \ref{definition-master})},
  \end{align*}
where $m_{s\tilde{s}} : s \to \tilde{s}$ is the unique
  1-cell in $\mathrm{Ind}(S)$. 
\end{proposition}

\begin{proof}

  \medskip
  \noindent\textbf{Reduction of the data.}

  \medskip
  We apply definition \ref{definition-master} with $m = \tilde{m} = m_{s\tilde{s}}$,
  $A = \tilde{A} = s$, and $B = \tilde{B} = \tilde{s}$.

  \smallskip
  \noindent\emph{Choice of auxiliary 1-cells.}
  The 1-cells $u_1 : \tilde{s} \to \tilde{s}$, $u_2 : s \to s$,
  $v_1 : \tilde{s} \to \tilde{s}$, $v_2 : s \to s$ required by
  definition~\ref{definition-master} are \emph{uniquely determined} in $\mathrm{Ind}(S)$:
  \[
    u_1 = m_{\tilde{s}\tilde{s}}, \quad
    u_2 = m_{ss}, \quad
    v_1 = m_{\tilde{s}\tilde{s}}, \quad
    v_2 = m_{ss}.
  \]
  There is no freedom of choice; any pair of objects in $\mathrm{Ind}(S)$ admits
  exactly one 1-cell between them.

  \smallskip
  \noindent\emph{Computation of the composite 1-cells in $\mathrm{Ind}(S')$.}
  Applying $\sigma = \tau_1 = \tau_2$ gives
  \[
    \sigma(u_1) = m'_{f(\tilde{s})f(\tilde{s})}, \quad
    \sigma(u_2) = m'_{f(s)f(s)}, \quad
    \sigma(m_{s\tilde{s}}) = m'_{f(s)f(\tilde{s})}.
  \]
  Since $\mathrm{Ind}(S')$ is indiscrete, the composite
  \[
    \sigma(u_1) \circ \sigma(m_{s\tilde{s}}) \circ \sigma(u_2)
    \;=\; m'_{f(\tilde{s})f(\tilde{s})} \circ
    m'_{f(s)f(\tilde{s})} \circ
    m'_{f(s)f(s)}
    \;=\; m'_{f(s)f(\tilde{s})}
    \;=\; \sigma(m_{s\tilde{s}}).
  \]
  By the same computation with $v_1, v_2$ (which equal $u_1, u_2$), the
  composites for $\Psi$ and $\widetilde{\Psi}$ likewise both equal
  $m'_{f(s)f(\tilde{s})}$.

  \smallskip
  \noindent\emph{The required 2-morphisms.}
  All four 2-morphisms $\Phi, \widetilde{\Phi}, \Psi, \widetilde{\Psi}$
  reduce to endomorphisms of the single 1-cell
  $m'_{f(s)f(\tilde{s})}$:
  \begin{align}
    \Phi,\,\widetilde{\Phi},\,\Psi,\,\widetilde{\Psi}
    \;:\;
    m'_{f(s)f(\tilde{s})} \Rightarrow m'_{f(s)f(\tilde{s})}
    \quad \text{in } \mathrm{Ind}(S').
    \label{eq:2cells}
  \end{align}
  The unique 2-cell in $\mathrm{Ind}(S')(f(s),\,f(\tilde{s}))$ is labeled
  $\tau'(f(s), f(\tilde{s})) \in \{\top, \bot\}$. By our convention, this
  constitutes a genuine 2-morphism if and only if
  $\tau'(f(s), f(\tilde{s})) = \top$, equivalently $f(s) \simeq f(\tilde{s})$
  in $S'$.

  \begin{itemize}
  
  \item \textbf{Proof of $(\Leftarrow)$: $f(s) \simeq f(\tilde{s})$
    implies $m_{s\tilde{s}} \simeq m_{s\tilde{s}}$.} \\
  Suppose $f(s) \simeq f(\tilde{s})$ in $S'$, so $\tau'(f(s),f(\tilde{s}))
    = \top$. Choose the auxiliary 1-cells as above. By~\eqref{eq:2cells}, each
  of the four required 2-morphisms is the unique 2-cell
  \[
    \Phi \;=\; \widetilde{\Phi} \;=\; \Psi \;=\; \widetilde{\Psi}
    \;:=\; \top
    \;\in\; \mathrm{Ind}(S')(f(s),\,f(\tilde{s})).
  \]
  Since $\tau'(f(s), f(\tilde{s})) = \top$, this is a genuine 2-morphism
  in $\mathrm{Ind}(S')$. Hence all four 2-morphisms of definition~\ref{definition-master}
  exist, and $m_{s\tilde{s}} \simeq m_{s\tilde{s}}$.

 \item \textbf{Proof of $(\Rightarrow)$: $m_{s\tilde{s}} \simeq m_{s\tilde{s}}$ implies
    $f(s) \simeq f(\tilde{s})$.}

  \medskip
  Suppose $m_{s\tilde{s}} \simeq m_{s\tilde{s}}$, so there exist 1-cells and 2-morphisms as in
  definition~\ref{definition-master}. As shown in the reduction step, the auxiliary
  1-cells are uniquely forced, and the four 2-morphisms all take the
  form~\eqref{eq:2cells}. In particular, $\Phi$ is a genuine 2-morphism in
  $\mathrm{Ind}(S')(f(s),\,f(\tilde{s}))$. The existence of any genuine 2-morphism in
  that hom-category requires $\tau'(f(s), f(\tilde{s})) = \top$, i.e.\
  $f(s) \simeq f(\tilde{s})$ in $S'$.
  \end{itemize}
  Both implications having been established, the proof is complete.
\end{proof}

\end{document}
